% ============================================================================
%  PRECÁLCULO — Sesión 4
%  Semana 4: 31 de agosto al 5 de septiembre de 2026
%  1.4 Conjuntos
%  Duración prevista: 3 horas
% ============================================================================
\documentclass[aspectratio=169,10pt]{beamer}

\usepackage[utf8]{inputenc}
\IfFileExists{spanish.ldf}{\usepackage[spanish,es-noshorthands]{babel}}{}
\usepackage{amsmath,amssymb}
\usepackage{booktabs}
\usepackage{tikz}

\usetheme{Uniminuto}

\uniasignatura{Precálculo}
\uniprofesor{Andrés Felipe Puerta Velez}
\unicorreo{andres.puerta-v@uniminuto.edu.co}
\uninrc{MATE1161}
\unisemana{Sesión 4 — Semana 4 · 31 de agosto al 5 de septiembre de 2026}

\newcommand{\listatiny}{\setbeamerfont{itemize/enumerate body}{size=\tiny}\tiny}

% --- diagrama de Venn de dos conjuntos -------------------------------------
% \vennII{etiqueta A}{etiqueta B}{sombreado: 0 nada,1 union,2 interseccion,
%         3 A-B, 4 complemento de la union}
\newcommand{\vennII}[3]{%
  \begin{tikzpicture}[scale=0.45]
    \draw[uniGris,fill=uniGris!6] (-2.6,-1.9) rectangle (4.6,1.9);
    \node[font=\tiny,color=uniGris] at (4.2,1.6) {$U$};
    \begin{scope}
      \clip (0,0) circle (1.5);
      \ifnum#3=2 \fill[uniNaranja!45] (2,0) circle (1.5); \fi
    \end{scope}
    \ifnum#3=1
      \fill[uniAzul!25] (0,0) circle (1.5);
      \fill[uniAzul!25] (2,0) circle (1.5);
    \fi
    \ifnum#3=3
      \begin{scope}
        \clip (0,0) circle (1.5);
        \fill[uniAzul!25] (-2.6,-1.9) rectangle (4.6,1.9);
        \fill[white] (2,0) circle (1.5);
      \end{scope}
    \fi
    \ifnum#3=4
      \fill[uniTitulo!25] (-2.6,-1.9) rectangle (4.6,1.9);
      \fill[white] (0,0) circle (1.5);
      \fill[white] (2,0) circle (1.5);
    \fi
    \draw[uniAzul,thick] (0,0) circle (1.5);
    \draw[uniTitulo,thick] (2,0) circle (1.5);
    \node[font=\tiny,color=uniAzul] at (-1.0,1.2) {#1};
    \node[font=\tiny,color=uniTitulo] at (3.0,1.2) {#2};
  \end{tikzpicture}}

\begin{document}

\portada

\begin{frame}{Orden del día}
  \begin{enumerate}
    \item Llamado a lista y revisión de los ejercicios propuestos
    \item Repaso: dónde quedamos
    \item \azul{Tema 1.4: notación, pertenencia, subconjuntos y cardinal}
    \item Operaciones: unión, intersección, diferencia y complemento
    \item Leyes de conjuntos y De Morgan
    \item Conteo con diagramas de Venn, taller y cierre
  \end{enumerate}

  \vspace{0.25cm}
  {\footnotesize\color{uniGris} Sesión de 3 horas. Este es el tema que más
  directamente reaparece en su carrera: \destacar{estructuras discretas,
  lógica, bases de datos y filtros de consulta} son conjuntos con otro
  nombre.}
\end{frame}

\begin{frame}{Repaso: dónde quedamos}
  \footnotesize
  En la sesión 3 operamos con racionales: simplificar, mcm, producto,
  cociente y fracciones complejas.

  \vspace{0.2cm}
  \begin{alertblock}{La pregunta de hoy}
    \footnotesize
    Hasta ahora trabajamos con \azul{números sueltos}. Hoy trabajamos con
    \destacar{colecciones de objetos} y con las operaciones que se hacen entre
    ellas. Ya usamos conjuntos sin nombrarlos: $\mathbb{N}$, $\mathbb{Z}$,
    $\mathbb{Q}$ y $\mathbb{R}$ de la sesión 1 \azul{son conjuntos}, y
    $\mathbb{R} = \mathbb{Q}\cup\mathbb{I}$ es una unión.
  \end{alertblock}

  \vspace{0.2cm}
  \begin{block}{Dónde lo va a volver a ver}
    \footnotesize
    Un \azul{\texttt{WHERE}} en una consulta es una intersección; un
    \azul{\texttt{OR}} es una unión; un \azul{\texttt{NOT}} es un complemento.
    Lo de hoy es la matemática de eso.
  \end{block}
\end{frame}

\seccion{1.4 Conjuntos}

\begin{frame}{Notación y pertenencia}
  \scriptsize
  \begin{columns}[T,onlytextwidth]
    \begin{column}{0.48\textwidth}
      \begin{block}{Por extensión}
        \tiny
        Se listan los elementos:
        \[ A = \{1,\,2,\,3,\,4,\,5\} \]
        \azul{No se repiten} elementos y \azul{no importa el orden}:
        $\{1,2\} = \{2,1\}$.
      \end{block}
    \end{column}
    \begin{column}{0.48\textwidth}
      \begin{block}{Por comprensión}
        \tiny
        Se da la propiedad que los define:
        \[ A = \{x \in \mathbb{N} \;:\; 1 \le x \le 5\} \]
        Se lee: ``los $x$ naturales \azul{tales que} están entre 1 y 5''.
      \end{block}
    \end{column}
  \end{columns}

  \vspace{0.12cm}
  \begin{center}
  \begin{tabular}{@{}p{2.2cm}p{4.0cm}p{5.4cm}@{}}
    \toprule
    \textbf{Símbolo} & \textbf{Se lee} & \textbf{Ejemplo} \\
    \midrule
    $3 \in A$ & 3 pertenece a $A$ & Relaciona \azul{elemento} con conjunto \\
    $8 \notin A$ & 8 no pertenece a $A$ & \\
    $B \subseteq A$ & $B$ está contenido en $A$ &
      Relaciona \azul{conjunto} con conjunto \\
    $|A|$ o $n(A)$ & Cardinal de $A$ & $|A| = 5$ \\
    $\varnothing$ & Conjunto vacío & $|\varnothing| = 0$ \\
    \bottomrule
  \end{tabular}
  \end{center}

  \vspace{0.05cm}
  {\tiny\color{uniGris} \destacar{Error clásico:} escribir $\{3\} \in A$. Lo
  correcto es $3 \in A$ o bien $\{3\} \subseteq A$. El símbolo $\in$ es para
  elementos; $\subseteq$, para conjuntos.}
\end{frame}

\begin{frame}{Subconjuntos y conjunto de partes}
  \scriptsize
  \begin{columns}[T,onlytextwidth]
    \begin{column}{0.48\textwidth}
      \begin{block}{Definiciones}
        \tiny
        $B \subseteq A$ si \azul{todo} elemento de $B$ está en $A$.

        \vspace{0.08cm}
        $B \subset A$ (propio) si además $B \neq A$.

        \vspace{0.08cm}
        $A = B$ si $A \subseteq B$ y $B \subseteq A$.

        \vspace{0.08cm}
        \destacar{$\varnothing \subseteq A$ siempre}, para cualquier $A$.
      \end{block}
    \end{column}
    \begin{column}{0.48\textwidth}
      \begin{exampleblock}{Conjunto de partes}
        \tiny
        $\mathcal{P}(A)$ es el conjunto de \azul{todos} los subconjuntos de
        $A$. Si $A = \{a,b,c\}$:
        \[ \varnothing,\ \{a\},\ \{b\},\ \{c\},\ \{a,b\},\ \{a,c\},\
           \{b,c\},\ \{a,b,c\} \]
        \[ |\mathcal{P}(A)| = 2^{|A|} = 2^3 = \destacar{8} \]
      \end{exampleblock}
    \end{column}
  \end{columns}

  \vspace{0.1cm}
  \begin{alertblock}{Por qué $2^n$, y por qué le va a servir}
    \scriptsize
    Cada elemento tiene dos opciones: \azul{entra o no entra} en el
    subconjunto. Con $n$ elementos hay $2\cdot 2\cdots 2 = 2^n$ combinaciones.
    \destacar{Es el mismo razonamiento que cuenta cuántos números representa un
    byte:} 8 bits, cada uno con dos estados, $2^8 = 256$.
  \end{alertblock}
\end{frame}

\begin{frame}{Las cuatro operaciones}
  \scriptsize
  \begin{center}
  \begin{tabular}{@{}cccc@{}}
    \vennII{$A$}{$B$}{1} & \vennII{$A$}{$B$}{2} &
    \vennII{$A$}{$B$}{3} & \vennII{$A$}{$B$}{4} \\[-0.15em]
    {\tiny $A \cup B$} & {\tiny $A \cap B$} &
    {\tiny $A - B$} & {\tiny $(A \cup B)^{c}$} \\
  \end{tabular}
  \end{center}

  \vspace{0.1cm}
  \begin{center}
  \begin{tabular}{@{}p{2.6cm}p{4.2cm}p{4.8cm}@{}}
    \toprule
    \textbf{Operación} & \textbf{Definición} & \textbf{Se lee} \\
    \midrule
    $A \cup B$ & $\{x : x\in A$ \azul{o} $x \in B\}$ & Todo lo que esté en
      alguno de los dos \\
    $A \cap B$ & $\{x : x\in A$ \azul{y} $x \in B\}$ & Solo lo que esté en los
      dos \\
    $A - B$ & $\{x : x\in A$ y $x \notin B\}$ & Lo de $A$ quitando lo de $B$ \\
    $A^{c}$ & $\{x \in U : x \notin A\}$ & Todo lo que \azul{no} está en $A$ \\
    \bottomrule
  \end{tabular}
  \end{center}

  \vspace{0.05cm}
  {\tiny\color{uniGris} El complemento \destacar{siempre} se toma respecto de
  un conjunto universal $U$: sin decir cuál es $U$, la expresión $A^c$ no
  significa nada.}
\end{frame}

\begin{frame}{Ejemplo guiado — todas las operaciones con los mismos datos}
  \scriptsize
  \[ U = \{1,2,3,4,5,6,7,8,9,10\} \qquad A = \{1,2,3,4,5\}
     \qquad B = \{4,5,6,7\} \]

  \vspace{0.1cm}
  \begin{center}
  \begin{tabular}{@{}lll@{}}
    \toprule
    \textbf{Operación} & \textbf{Resultado} & \textbf{Cómo se obtuvo} \\
    \midrule
    $A \cup B$ & $\{1,2,3,4,5,6,7\}$ & Todo, sin repetir el 4 y el 5 \\
    $A \cap B$ & $\{4,5\}$ & Lo que aparece en las dos listas \\
    $A - B$ & $\{1,2,3\}$ & De $A$, quitando 4 y 5 \\
    $B - A$ & $\{6,7\}$ & \destacar{No es lo mismo que $A-B$} \\
    $A^{c}$ & $\{6,7,8,9,10\}$ & Lo de $U$ que no está en $A$ \\
    $A \triangle B$ & $\{1,2,3,6,7\}$ & Lo que está en uno pero no en ambos \\
    \bottomrule
  \end{tabular}
  \end{center}

  \vspace{0.08cm}
  \begin{alertblock}{Dos observaciones}
    \tiny
    La \destacar{unión y la intersección son conmutativas}, pero la
    \destacar{diferencia no}: $A-B \neq B-A$. \quad
    La diferencia simétrica cumple $A \triangle B = (A\cup B) - (A\cap B)$:
    es ``lo que no comparten''.
  \end{alertblock}
\end{frame}

\begin{frame}{Leyes de conjuntos y De Morgan}
  \scriptsize
  \begin{center}
  \begin{tabular}{@{}p{2.6cm}p{4.4cm}p{4.6cm}@{}}
    \toprule
    \textbf{Ley} & \textbf{Con $\cup$} & \textbf{Con $\cap$} \\
    \midrule
    Conmutativa & $A\cup B = B\cup A$ & $A\cap B = B\cap A$ \\
    Asociativa & $(A\cup B)\cup C = A\cup(B\cup C)$ &
      $(A\cap B)\cap C = A\cap(B\cap C)$ \\
    Distributiva & $A\cap(B\cup C) = (A\cap B)\cup(A\cap C)$ &
      $A\cup(B\cap C) = (A\cup B)\cap(A\cup C)$ \\
    Identidad & $A\cup\varnothing = A$ & $A\cap U = A$ \\
    Complemento & $A\cup A^{c} = U$ & $A\cap A^{c} = \varnothing$ \\
    \bottomrule
  \end{tabular}
  \end{center}

  \vspace{0.1cm}
  \begin{exampleblock}{Leyes de De Morgan}
    \scriptsize
    \[ (A\cup B)^{c} = A^{c}\cap B^{c} \qquad\qquad
       (A\cap B)^{c} = A^{c}\cup B^{c} \]
    Verificación con los datos anteriores:
    $(A\cup B)^c = \{8,9,10\}$ \ y \ $A^c\cap B^c
    = \{6,7,8,9,10\}\cap\{1,2,3,8,9,10\} = \{8,9,10\}$ \checkmark
  \end{exampleblock}

  \vspace{0.05cm}
  {\tiny\color{uniGris} En palabras: \destacar{negar un ``o'' lo convierte en
  ``y''}, y al revés. Es exactamente lo que hace \azul{\texttt{!(a || b)} =
  \texttt{!a \&\& !b}} en cualquier lenguaje de programación.}
\end{frame}

\begin{frame}{Ejercicio 1 — en clase}
  \begin{alertblock}{Individual (12 minutos)}
    \scriptsize
    Sean \ $U = \{x \in \mathbb{N} : 1 \le x \le 12\}$, \quad
    $A = \{2,4,6,8,10\}$, \quad $B = \{1,2,3,4,5\}$, \quad
    $C = \{5,6,7,8\}$.

    \vspace{0.25cm}
    Determine, listando los elementos:

    \vspace{0.15cm}
    \begin{columns}[T,onlytextwidth]
      \begin{column}{0.48\textwidth}
        \begin{enumerate}\setlength{\itemsep}{0.02cm}
          \item $A \cup B$
          \item $A \cap C$
          \item $B - C$
          \item $A^{c}$
          \item $(A \cap B)^{c}$
        \end{enumerate}
      \end{column}
      \begin{column}{0.48\textwidth}
        \begin{enumerate}\setlength{\itemsep}{0.02cm}
          \setcounter{enumi}{5}
          \item $A \cap (B \cup C)$
          \item $(A\cap B)\cup(A\cap C)$
          \item $B \triangle C$
          \item $|A \cup B \cup C|$
          \item $|\mathcal{P}(C)|$
        \end{enumerate}
      \end{column}
    \end{columns}
  \end{alertblock}
\end{frame}

\begin{frame}{Ejercicio 1 — solución}
  \tiny
  \begin{columns}[T,onlytextwidth]
    \begin{column}{0.49\textwidth}
      \begin{block}{Puntos 1 a 5}
        \tiny
        \begin{enumerate}\setlength{\itemsep}{-0.02cm}
          \item $\{1,2,3,4,5,6,8,10\}$
          \item $\{6,8\}$
          \item $\{1,2,3,4\}$
          \item $\{1,3,5,7,9,11,12\}$
          \item $A\cap B = \{2,4\}$, luego
                $\{1,3,5,6,7,8,9,10,11,12\}$
        \end{enumerate}
      \end{block}
    \end{column}
    \begin{column}{0.49\textwidth}
      \begin{block}{Puntos 6 a 10}
        \tiny
        \begin{enumerate}\setlength{\itemsep}{-0.02cm}
          \setcounter{enumi}{5}
          \item $B\cup C = \{1,\dots,8\}$, luego $\{2,4,6,8\}$
          \item $\{2,4\}\cup\{6,8\} = \{2,4,6,8\}$
          \item $\{1,2,3,4,6,7,8\}$
          \item $|\{1,\dots,8\}\cup\{10\}| = 9$
          \item $2^{4} = 16$
        \end{enumerate}
      \end{block}
    \end{column}
  \end{columns}

  \vspace{0.1cm}
  \begin{alertblock}{Compare el 6 con el 7}
    \tiny
    Dan lo mismo, y no es casualidad: es la \destacar{ley distributiva}
    $A\cap(B\cup C) = (A\cap B)\cup(A\cap C)$ verificada con números. Cuando
    una ley se cumple en un ejemplo concreto, deja de ser una fórmula que
    memorizar.
  \end{alertblock}
\end{frame}

\begin{frame}{Conteo con diagramas de Venn}
  \scriptsize
  \begin{columns}[T,onlytextwidth]
    \begin{column}{0.44\textwidth}
      \centering
      \begin{tikzpicture}[scale=0.72]
        \draw[uniGris,fill=uniGris!6] (-2.4,-2.0) rectangle (4.6,2.0);
        \node[font=\tiny,color=uniGris] at (4.2,1.7) {$U = 100$};
        \draw[uniAzul,thick,fill=uniAzul,fill opacity=0.18] (0,0) circle (1.6);
        \draw[uniTitulo,thick,fill=uniTitulo,fill opacity=0.18] (2.1,0) circle (1.6);
        \node[font=\tiny,color=uniAzul] at (-1.1,1.3) {Python};
        \node[font=\tiny,color=uniTitulo] at (3.2,1.3) {Java};
        \node[font=\scriptsize] at (-0.5,0) {35};
        \node[font=\scriptsize] at (1.05,0) {25};
        \node[font=\scriptsize] at (2.6,0) {20};
        \node[font=\scriptsize,color=uniGris] at (3.9,-1.6) {20};
      \end{tikzpicture}
    \end{column}
    \begin{column}{0.54\textwidth}
      \begin{block}{El problema}
        \tiny
        De 100 estudiantes: \azul{60} programan en Python, \azul{45} en Java y
        \azul{25} en ambos.
        \vspace{0.06cm}
        \[ |P\cup J| = |P| + |J| - |P\cap J| = 60+45-25 = \destacar{80} \]
        Solo Python: $60-25 = 35$. \quad Solo Java: $45-25 = 20$. \quad
        Ninguno: $100-80 = \destacar{20}$.
      \end{block}
    \end{column}
  \end{columns}

  \vspace{0.08cm}
  \begin{alertblock}{Principio de inclusión-exclusión}
    \tiny
    \[ |A\cup B| = |A|+|B|-|A\cap B| \]
    \[ |A\cup B\cup C| = |A|+|B|+|C|-|A\cap B|-|A\cap C|-|B\cap C|+|A\cap B\cap C| \]
    Se resta lo que se contó dos veces y se devuelve lo que se restó de más.
    \destacar{Siempre se empieza a llenar el diagrama por el centro.}
  \end{alertblock}
\end{frame}

\begin{frame}{Ejercicio 2 — en clase (taller principal)}
  \begin{alertblock}{En grupos de 3 (25 minutos)}
    \scriptsize
    En un grupo de \destacar{120 estudiantes} de Ingeniería de Software:

    \vspace{0.15cm}
    \centering
    \begin{tabular}{@{}lr@{\hspace{1.4cm}}lr@{}}
      Saben Python & 70 & Python y Java & 25 \\
      Saben Java   & 50 & Python y C++  & 20 \\
      Saben C++    & 40 & Java y C++    & 15 \\
      & & Los tres lenguajes & 10 \\
    \end{tabular}

    \vspace{0.2cm}
    \raggedright
    \begin{enumerate}\setlength{\itemsep}{-0.05cm}
      \item Dibuje el diagrama de Venn y complete \destacar{las siete
            regiones}, empezando por el centro.
      \item ¿Cuántos saben al menos un lenguaje?
      \item ¿Cuántos no saben ninguno?
      \item ¿Cuántos saben \destacar{exactamente uno}?
      \item ¿Cuántos saben \destacar{exactamente dos}?
      \item Verifique que las siete regiones sumen lo del punto 2.
    \end{enumerate}
  \end{alertblock}
\end{frame}

\begin{frame}{Ejercicio 2 — solución}
  \scriptsize
  \begin{columns}[T,onlytextwidth]
    \begin{column}{0.44\textwidth}
      \centering
      \begin{tikzpicture}[scale=0.68]
        \draw[uniGris,fill=uniGris!6] (-2.6,-2.9) rectangle (4.6,2.6);
        \node[font=\tiny,color=uniGris] at (4.0,2.3) {$U=120$};
        \draw[uniAzul,thick] (0,0.5) circle (1.7);
        \draw[uniTitulo,thick] (2.0,0.5) circle (1.7);
        \draw[uniNaranja,thick] (1.0,-1.1) circle (1.7);
        \node[font=\tiny,color=uniAzul] at (-1.3,1.9) {Python};
        \node[font=\tiny,color=uniTitulo] at (3.3,1.9) {Java};
        \node[font=\tiny,color=uniNaranja] at (1.0,-2.6) {C++};
        \node[font=\scriptsize] at (-0.5,0.9) {35};
        \node[font=\scriptsize] at (2.5,0.9) {20};
        \node[font=\scriptsize] at (1.0,-1.9) {15};
        \node[font=\scriptsize] at (1.0,1.1) {15};
        \node[font=\scriptsize] at (0.1,-0.7) {10};
        \node[font=\scriptsize] at (1.9,-0.7) {5};
        \node[font=\scriptsize] at (1.0,-0.1) {10};
        \node[font=\scriptsize,color=uniGris] at (3.9,-2.5) {10};
      \end{tikzpicture}
    \end{column}
    \begin{column}{0.54\textwidth}
      \begin{block}{Los cálculos}
        \tiny
        Del centro hacia afuera:
        \[ P\cap J \text{ solo} = 25-10 = 15 \qquad
           P\cap C \text{ solo} = 20-10 = 10 \]
        \[ J\cap C \text{ solo} = 15-10 = 5 \]
        \[ \text{Solo } P = 70-15-10-10 = 35 \]
        \[ \text{Solo } J = 50-15-5-10 = 20 \qquad
           \text{Solo } C = 40-10-5-10 = 15 \]
      \end{block}
    \end{column}
  \end{columns}

  \vspace{0.05cm}
  \begin{alertblock}{Respuestas}
    \tiny
    \textbf{2.} $70+50+40-25-20-15+10 = \destacar{110}$. \quad
    \textbf{3.} $120-110 = \destacar{10}$. \quad
    \textbf{4.} $35+20+15 = \destacar{70}$. \quad
    \textbf{5.} $15+10+5 = \destacar{30}$. \quad
    \textbf{6.} $70+30+10 = 110$ \checkmark
  \end{alertblock}
\end{frame}

\begin{frame}{Ejercicio 3 — en clase}
  \begin{alertblock}{En parejas (12 minutos)}
    \scriptsize
    \begin{enumerate}\setlength{\itemsep}{0.08cm}
      \item Verifique con $A=\{1,2,3,4,5\}$, $B=\{4,5,6,7\}$ y
            $U=\{1,\dots,10\}$ que $(A\cap B)^{c} = A^{c}\cup B^{c}$.
      \item Escriba por comprensión: $\{2,4,6,8,10,\dots\}$ y
            $\{1,4,9,16,25\}$.
      \item ¿Verdadero o falso? Justifique.
            \begin{itemize}\setlength{\itemsep}{-0.04cm}
              \item[a)] $\varnothing \subseteq \{1,2\}$
              \item[b)] $\varnothing \in \{1,2\}$
              \item[c)] $\{1\} \in \{1,2\}$
              \item[d)] $\{1\} \subseteq \{1,2\}$
            \end{itemize}
      \item Si $|A| = 6$, ¿cuántos subconjuntos tiene $A$? ¿Cuántos son
            propios?
      \item Dé un ejemplo de conjuntos $A$ y $B$ con $A - B = B - A$.
    \end{enumerate}
  \end{alertblock}
\end{frame}

\begin{frame}{Ejercicio 3 — solución}
  \scriptsize
  \begin{columns}[T,onlytextwidth]
    \begin{column}{0.49\textwidth}
      \begin{block}{Puntos 1 y 2}
        \tiny
        \textbf{1.} $A\cap B = \{4,5\}$, luego
        $(A\cap B)^c = \{1,2,3,6,7,8,9,10\}$.

        $A^c = \{6,\dots,10\}$, \ $B^c = \{1,2,3,8,9,10\}$;
        su unión es $\{1,2,3,6,7,8,9,10\}$ \checkmark

        \vspace{0.1cm}
        \textbf{2.} $\{x\in\mathbb{N} : x \text{ es par},\ x\ge 2\}$ \ y \
        $\{n^2 : n\in\mathbb{N},\ 1\le n\le 5\}$
      \end{block}
    \end{column}
    \begin{column}{0.49\textwidth}
      \begin{alertblock}{Punto 3: el corazón del tema}
        \tiny
        \textbf{a) V:} el vacío está contenido en todo conjunto. \\
        \textbf{b) F:} $\varnothing$ no es \azul{elemento} de $\{1,2\}$; los
        elementos son 1 y 2. \\
        \textbf{c) F:} $\{1\}$ es un conjunto, no un elemento de $\{1,2\}$. \\
        \textbf{d) V:} $\{1\}$ sí está \azul{contenido}.
      \end{alertblock}
    \end{column}
  \end{columns}

  \vspace{0.05cm}
  \begin{exampleblock}{Puntos 4 y 5}
    \tiny
    \textbf{4.} $2^6 = \destacar{64}$ subconjuntos; propios,
    $64 - 1 = \destacar{63}$ (todos menos el mismo $A$).

    \vspace{0.06cm}
    \textbf{5.} Solo ocurre cuando $A = B$: entonces ambas diferencias son
    $\varnothing$. \destacar{Es el único caso}, y encontrarlo es el punto del
    ejercicio.
  \end{exampleblock}
\end{frame}

\begin{frame}{Errores frecuentes en esta sesión}
  \scriptsize
  \begin{center}
  \begin{tabular}{@{}p{5.2cm}p{6.6cm}@{}}
    \toprule
    \textbf{Error} & \textbf{Corrección} \\
    \midrule
    Escribir $\{3\} \in A$ &
      $\in$ es para elementos: $3 \in A$ o $\{3\} \subseteq A$ \\
    \addlinespace
    Confundir $\varnothing$ con $\{\varnothing\}$ &
      El primero es vacío; el segundo tiene \azul{un} elemento \\
    \addlinespace
    Suponer que $A-B = B-A$ &
      La diferencia \azul{no} es conmutativa \\
    \addlinespace
    Calcular $A^{c}$ sin definir $U$ &
      El complemento siempre es respecto de un universal \\
    \addlinespace
    En Venn: poner $|A\cap B|$ solo en la zona central &
      Ese número incluye a los que están en los tres; se resta primero \\
    \addlinespace
    Sumar $|A|+|B|$ y olvidar restar la intersección &
      Los comunes quedarían contados dos veces \\
    \bottomrule
  \end{tabular}
  \end{center}
\end{frame}

\begin{frame}{Cierre de la sesión}
  \begin{block}{Lo que debe quedar claro hoy}
    \scriptsize
    \begin{itemize}\setlength{\itemsep}{0.02cm}
      \item Un conjunto se define por \azul{extensión} o por
            \azul{comprensión}; no hay repeticiones ni orden.
      \item $\in$ relaciona elemento con conjunto; $\subseteq$, conjunto con
            conjunto.
      \item $|\mathcal{P}(A)| = 2^{|A|}$, por el argumento de ``entra o no
            entra''.
      \item Unión, intersección, diferencia y complemento, con sus diagramas.
      \item \azul{De Morgan}: negar un ``o'' lo vuelve ``y'', y viceversa.
      \item Inclusión-exclusión: se resta lo contado dos veces y se devuelve lo
            restado de más.
    \end{itemize}
  \end{block}

  \vspace{0.08cm}
  \begin{alertblock}{Trabajo independiente}
    \tiny
    Resolver los 20 ejercicios de la lámina siguiente. En la sesión 5 cerramos
    la unidad con \destacar{razones y proporciones}, y hacemos el repaso
    integrador y el simulacro del Examen 1.
  \end{alertblock}
\end{frame}

\begin{frame}{Ejercicios propuestos}
  \vspace{-0.1cm}
  \begin{beamercolorbox}[wd=\textwidth,sep=0.10cm]{block body}
    \centering\scriptsize
    Con $U=\{1,\dots,10\}$, \ $A=\{1,3,5,7,9\}$, \ $B=\{2,3,5,7\}$, \
    $C=\{1,2,3,4\}$.
  \end{beamercolorbox}
  \listatiny
  \vspace{0.15cm}
  \begin{columns}[T,onlytextwidth]
    \begin{column}{0.47\textwidth}
      \begin{enumerate}\setlength{\itemsep}{0.16cm}
        \item $A \cup B$
        \item $A \cap B$
        \item $A - C$
        \item $C - A$
        \item $B^{c}$
        \item $(A\cup C)^{c}$
        \item $A^{c}\cap C^{c}$
        \item $A \cap (B\cup C)$
        \item $B \triangle C$
        \item $|A\cup B\cup C|$
      \end{enumerate}
    \end{column}
    \begin{column}{0.50\textwidth}
      \begin{enumerate}\setlength{\itemsep}{0.16cm}
        \setcounter{enumi}{10}
        \item $|\mathcal{P}(B)|$
        \item ¿Es $B \subseteq A$?
        \item ¿Es $\{3,5\} \subseteq A\cap B$?
        \item Escriba $A$ por comprensión
        \item ¿Cuántos subconjuntos propios tiene $C$?
        \item De 80 personas, 50 leen X, 35 leen Y y 20 leen ambos:
              ¿cuántas no leen ninguno?
        \item En el caso anterior, ¿cuántas leen solo X?
        \item Enuncie las dos leyes de De Morgan
        \item ¿Qué operación de conjuntos corresponde a un \texttt{AND}?
        \item ¿Y a un \texttt{NOT}?
      \end{enumerate}
    \end{column}
  \end{columns}
\end{frame}

\begin{frame}{Respuestas}{Autoevaluación de los ejercicios propuestos}
  \listatiny
  \vspace{0.15cm}
  \begin{columns}[T,onlytextwidth]
    \begin{column}{0.47\textwidth}
      \begin{enumerate}\setlength{\itemsep}{0.16cm}
        \item $\{1,2,3,5,7,9\}$
        \item $\{3,5,7\}$
        \item $\{5,7,9\}$
        \item $\{2,4\}$
        \item $\{1,4,6,8,9,10\}$
        \item $\{6,8,10\}$
        \item $\{6,8,10\}$ \ (De Morgan: igual al 6)
        \item $\{1,3,5,7\}$
        \item $\{1,4,5,7\}$
        \item $|\{1,2,3,4,5,7,9\}| = 7$
      \end{enumerate}
    \end{column}
    \begin{column}{0.50\textwidth}
      \begin{enumerate}\setlength{\itemsep}{0.16cm}
        \setcounter{enumi}{10}
        \item $2^4 = 16$
        \item No: $2 \in B$ pero $2 \notin A$
        \item Sí
        \item $\{x\in\mathbb{N} : x \text{ impar},\ 1\le x\le 9\}$
        \item $2^4 - 1 = 15$
        \item $80 - (50+35-20) = 15$
        \item $50-20 = 30$
        \item $(A\cup B)^c = A^c\cap B^c$; \ $(A\cap B)^c = A^c\cup B^c$
        \item La intersección
        \item El complemento
      \end{enumerate}
    \end{column}
  \end{columns}
\end{frame}

\begin{frame}{Referencias}
  \footnotesize
  \begin{itemize}
    \item[{[1]}] Ramos del Olmo, Á. M. (2016). \textit{Matemáticas básicas
      para el acceso a la universidad}. Ediciones Pirámide.
      \texttt{https://elibro.net/es/ereader/uniminuto/49134}
    \item[{[2]}] González Pulido, J. W., Chacón Benavides, J. A. \&
      Fonseca Correa, L. Á. (2023). \textit{Matemática básica} (1.\textsuperscript{a} ed.).
      Universidad Mariana.
      \texttt{https://elibro.net/es/ereader/uniminuto/280647}
    \item[{[3]}] Placencia Valero, J. (2008). \textit{Compendio de matemática
      básica elemental} (1.\textsuperscript{a} ed.). Editorial Tébar Flores.
      \texttt{https://elibro.net/es/ereader/uniminuto/269758}
  \end{itemize}

  \vspace{0.35cm}
  \begin{center}
    {\Large\color{uniAzul}\textbf{¿Preguntas?}}
  \end{center}
\end{frame}

\end{document}
